% NEON-EPN: an ESWC Resources Track submission.
%
% Springer's llncs.cls is not redistributable, so this uses article.cls with
% geometry set to match. To submit: download llncs.cls and splncs04.bst from
% Springer, replace the preamble down to \begin{document} with
%   \documentclass[runningheads]{llncs}
% and switch \bibliographystyle{splncs04}.
%
% Prepared for dual-anonymous review: no author names, and the repository is
% referenced through an anonymised mirror. Replace both at camera-ready.

\documentclass[11pt]{article}
\usepackage[a4paper,margin=2.5cm]{geometry}
\usepackage[T1]{fontenc}
\usepackage[utf8]{inputenc}
\usepackage{booktabs}
\usepackage{amsmath}
\usepackage{graphicx}
\usepackage{listings}
\usepackage[hidelinks]{hyperref}
\usepackage{xcolor}

\lstset{basicstyle=\ttfamily\footnotesize, breaklines=true,
        frame=single, framesep=4pt, rulecolor=\color{gray!40},
        showstringspaces=false, columns=fullflexible}

\title{\textbf{NEON-EPN: A Semantic Layer over Continental Ecological
Monitoring, and the Temporal Graph Benchmark Derived from It}}
\author{Anonymous Author(s)\\\small Anonymous Institution}
\date{}

\begin{document}
\maketitle

\begin{abstract}
\noindent
Ecological observatories publish measurements as independent tabular products,
each with its own identifiers, sampling design and temporal resolution. The
scientific questions that motivate such observatories are integrative, but the
data are not. We present \textbf{NEON-EPN}, an OWL ontology and reproducible
pipeline that integrates five data products from the US National Ecological
Observatory Network with volunteer phenology observations from the USA National
Phenology Network into a single validated knowledge graph, and a derivation of
that graph into a Temporal Graph Benchmark dataset for link prediction. The
resource comprises 10.5 million triples over 39 sites and eleven years, passing
a gate of 22 SHACL shapes with zero violations, and exports 725{,}826
timestamped events across five relations with 40{,}842 held-out evaluation
positives. Two contributions go beyond the artifact. First, aligning both
networks to the Plant Phenology Ontology enables a cross-source query that would
otherwise require a bespoke crosswalk: at 27 co-located sites, volunteer and
professional observers report bud-break dates differing by a median of 1.6 days,
with disagreement tracking sample size rather than observer type. Second, the
derived benchmark contains four prediction relations whose optimal baselines
differ, including one on which memorization scores \emph{below} the random
floor --- a property we have not found in existing temporal graph benchmarks.
Ontology, mappings, shapes, pipeline, dataset and baselines are released under
CC~BY~4.0 and MIT.

\medskip
\noindent\textbf{Keywords:} knowledge graph $\cdot$ ontology $\cdot$ SHACL
$\cdot$ RML $\cdot$ ecological data integration $\cdot$ temporal graph
benchmark $\cdot$ data quality
\end{abstract}

%======================================================================
\section{Introduction}

The US National Ecological Observatory Network (NEON) operates 81 field sites
and publishes roughly 180 standardised data products spanning organisms,
sensors and airborne remote sensing. Its stated purpose is to enable
continental-scale ecological inference. Its published form is a catalogue of
independent tables.

This gap is not incidental. A question such as \emph{does host community
composition predict pathogen prevalence in ticks} requires joining small mammal
trapping, tick pathogen assays, drag-cloth sampling and climate telemetry ---
four products whose records share no common key, identify locations at
different granularities, and encode identity in incompatible ways. In practice
each research group writes the joins again, with assumptions that are neither
recorded nor reviewable.

We argue that this is precisely the problem a semantic layer addresses, and
that ecological observatories are an underexploited application domain for
knowledge graph technology. We present NEON-EPN, comprising:

\begin{enumerate}
  \item an OWL ontology (25 classes) reusing SOSA, GeoSPARQL, QUDT, OWL-Time,
        NCBITaxon, ENVO, the Plant Phenology Ontology and PROV-O, with the
        import-versus-reference decision made explicitly for each;
  \item declarative RML mappings and a reproducible HPC pipeline producing
        10{,}536{,}423 triples over 39 sites and eleven years;
  \item a SHACL gate of 22 node shapes including cross-record integrity
        constraints, passing with zero violations;
  \item a derivation of the graph into Temporal Graph Benchmark format, with
        five relations, 725{,}826 events and reference baselines; and
  \item an integration of a second, independently operated observation network
        through shared ontology terms, and a cross-source comparison that the
        integration makes possible.
\end{enumerate}

Section~\ref{sec:validation} reports something we consider as valuable as the
artifact: the SHACL gate caught three classes of referential defect that no
structural check on any single product would surface, and \emph{failed} to
catch four data-modelling errors that only distributional inspection revealed.
We report both, because a resource paper that presents only its successes
misrepresents what validation does.

%======================================================================
\section{Related Work}

\paragraph{Semantic layers over observation networks.}
The closest analogue is the Biodiversity Exploratories Semantic Units knowledge
graph, which applies semantic units and a reuse-heavy vocabulary stack to a
German long-term ecological research network. FooDS constructs an ontology and
knowledge graph over a forest observatory in Borneo using YARRRML mappings ---
the same declarative family we use, applied to wildlife sensor data. Both
demonstrate that the approach transfers across observatories; neither derives a
machine learning benchmark from the result, which is the step we add.

For NEON specifically, prior semantic work is confined to metadata. NEON
publishes EML descriptors, and a recent effort mapped units from NEON, EDI and
DataONE to QUDT. Biodiversity ontologies including BCO and ENVO cite NEON
specimens among their motivating use cases. To our knowledge no formal semantic
layer over NEON \emph{observation} data has been published.

\paragraph{Ontology engineering practice.}
Recent domain ontologies in environmental and exposure science --- OntoPFAS for
perfluoroalkyl substances is a close methodological analogue --- follow the
Linked Open Terms methodology: requirements elicited as competency questions,
glossaries merged, existing vocabularies reused, then formalisation and
evaluation against those same questions. We follow the same discipline, and
Section~\ref{sec:cq} states the questions the ontology was built to answer.
Where our resource differs is in what follows construction: OntoPFAS publishes
sample instances, whereas the contribution here includes a graph instantiated
at scale over eleven years and a machine learning benchmark derived from it.

\paragraph{Vocabulary reuse.}
SOSA/SSN is the established vocabulary for observations and is used at scale in
geospatial knowledge graphs. The Plant Phenology Ontology is directly relevant
here: it was designed to integrate USA-NPN and NEON phenology records, and
Section~\ref{sec:crosssource} reports what happens when that design intention
is exercised.

\paragraph{Temporal graph benchmarks.}
TGB~2.0 provides eight datasets spanning temporal knowledge graphs and
heterogeneous graphs, with a standardised one-vs-many evaluation protocol.
Recent work notes the scarcity of temporal heterogeneous benchmarks and the
resulting reliance on a small pool of datasets. Conversions between benchmark
formats have moved from temporal graph representations into relational ones; we
are not aware of prior work in the opposite direction, deriving a temporal
graph benchmark from RDF.

%======================================================================
\section{Ontology Design}
\label{sec:ontology}

The ontology (\texttt{neonepn}, v0.6.0, 25 own classes, 593 triples) is
organised around the observation as a first-class entity. This is not a
stylistic choice: RDF triples have three positions, all occupied, and no
natural place for a timestamp. Reifying each measurement as a
\texttt{sosa:Observation} node gives the timestamp somewhere to live, and is
what makes the temporal export of Section~\ref{sec:benchmark} possible at all.

\subsection{Reuse, and the decision not to import}

Vocabulary reuse is the primary evaluation criterion for a published ontology,
but \emph{how} to reuse is a design decision with consequences. We distinguish
two cases.

\textbf{SOSA is imported in full.} It is load-bearing: four superficially
unrelated event types --- an animal captured in a trap, a blood assay, a tick
tested for pathogens, a temperature reading --- all become SOSA subclasses
sharing one shape. That shared shape is what allows a single SPARQL query to
span pathogen assays and sensor telemetry.

\textbf{GeoSPARQL, QUDT, OWL-Time, NCBITaxon, ENVO and PPO are referenced by
IRI} in the MIREOT style, without import. Importing OWL-Time alone pulled in
several hundred axioms to support three terms, slowing reasoning for no
expressive gain. We declare the terms we use and link to the authoritative
definitions.

\begin{table}[t]
\centering
\caption{External vocabularies and the form of reuse.}
\label{tab:reuse}
\small
\begin{tabular}{llp{6.6cm}}
\toprule
\textbf{Vocabulary} & \textbf{Reuse} & \textbf{Role} \\
\midrule
SOSA/SSN      & imported  & observation, sampling, sensor, feature of interest \\
GeoSPARQL     & referenced & site and plot geometry as WKT \\
QUDT          & referenced & quantity values with explicit units \\
OWL-Time      & referenced & daily intervals for aggregated climate \\
NCBITaxon     & referenced & host, pathogen and vector taxa \\
ENVO          & referenced & environmental context \\
PPO           & referenced & phenophase alignment (Sec.~\ref{sec:crosssource}) \\
PROV-O        & referenced & contributing network as an agent \\
\bottomrule
\end{tabular}
\end{table}

\subsection{Modelling decisions that the data forced}

Three decisions are worth reporting because they were not apparent from the
product documentation.

\emph{Daily climate needs an interval, not an instant.} A daily mean
temperature is not observed at midnight. Aggregated readings carry
\texttt{sosa:phenomenonTime} pointing at a \texttt{time:Interval}, with
\texttt{sosa:resultTime} marking when the aggregate was computed.

\emph{Proximity is not containment.} USA-NPN stations are volunteer-chosen
points near, but not inside, NEON sites. They are modelled as
\texttt{ObservationSite} with \texttt{nearFieldSite} and a recorded distance,
deliberately \emph{not} \texttt{withinSpatialUnit}. Asserting containment would
license false inference through the transitive spatial hierarchy; a SHACL shape
enforces the prohibition rather than leaving it to documentation.

\emph{Detection in an organism and detection at a place are different
relations.} Host-borne assays attach a pathogen to an animal; tick assays have
no host animal and attach to a location. Collapsing them would merge two
prediction targets with different semantics --- and, as
Section~\ref{sec:benchmark} shows, opposite statistical behaviour.

\subsection{Competency questions}
\label{sec:cq}

The ontology was developed against a set of questions that a practitioner
should be able to put to the integrated graph and that no single NEON product
can answer. Table~\ref{tab:cq} lists them with the construct each exercises;
all are answerable over the released graph, and the SPARQL is distributed with
the resource.

\begin{table}[t]
\centering
\caption{Competency questions. Each requires at least two source products, and
CQ7 requires two independently operated observation networks.}
\label{tab:cq}
\small
\begin{tabular}{llp{5.4cm}}
\toprule
& \textbf{Question} & \textbf{Construct exercised} \\
\midrule
CQ1 & Which pathogens have been detected in which host species, and where?
    & \texttt{detectedIn} through \texttt{ofTaxon} and the spatial hierarchy \\
CQ2 & What is the prevalence of a pathogen at a site, counting assays that
      found nothing? & observed negatives via \texttt{hasAssayResult} \\
CQ3 & Which tick species were collected at a plot, and in which years?
    & \texttt{observedAt} with \texttt{sosa:resultTime} \\
CQ4 & What were mean temperature and humidity in the 14 days before a given
      assay, at that site? & sensor telemetry joined to biology by place and
      time \\
CQ5 & Do plots that tested positive for a pathogen test positive again later?
    & temporal self-join over \texttt{detectedAt} \\
CQ6 & On what day of year does a phenophase first appear at a site, and how
      does it vary between years? & \texttt{expressedAt} with status
      transitions \\
CQ7 & Do two observation networks report the same phenophase timing at the
      same location? & \texttt{observedByNetwork} over PPO-aligned terms
      (Sec.~\ref{sec:crosssource}) \\
\bottomrule
\end{tabular}
\end{table}

CQ2 is worth singling out. Prevalence is the ratio of detections to assays
performed, and most published representations record only detections --- the
denominator is lost. Modelling the assay as the observation and the result as
its outcome keeps both, which is what makes the 347{,}615 observed negatives of
Section~\ref{sec:benchmark} available rather than inferred.

CQ7 could not be answered before the second source was integrated, and is the
question the cross-source section reports on.

\subsection{Formal constraints}

Disjointness and defined classes are asserted where the domain supports them,
so that inconsistency is detectable by a reasoner rather than only by SHACL:

\begin{align*}
\mathsf{HostTaxon} \sqcap \mathsf{PathogenTaxon} \sqcap
  \mathsf{VectorTaxon} \sqcap \mathsf{PlantTaxon} &\sqsubseteq \bot \\
\mathsf{PathogenAssay} \sqcap \mathsf{DragSample} \sqcap
  \mathsf{PhenophaseObservation} &\sqsubseteq \bot \\
\mathsf{PhenophaseStatus} &\equiv \{\mathsf{Expressed},
  \mathsf{NotExpressed}, \mathsf{StatusUncertain}\} \\
\mathsf{VectorOccurrence} &\sqsubseteq\ =\!1\ \mathsf{ofVectorTaxon}
\end{align*}

\noindent
The closed status set matters in practice: it is what rejected an invalid
phenophase value during the cross-source ingestion, before any data reached the
export.

%======================================================================
\section{Construction Pipeline}

Ingestion is declarative. Each NEON product is normalised into staging CSVs and
materialised through RML mappings (99 TriplesMaps), so the mapping from source
column to triple is inspectable rather than embedded in procedural code. Joins
are resolved in staging because morph-kgc join conditions over large CSVs were
prohibitively slow.

The pipeline runs as Slurm jobs on a university HPC cluster in two phases, a
split forced by the API rather than chosen. A fully parallel 39-way download
failed against NEON's rate limiter even with the hourly quota untouched, so
acquisition is sequential while transformation --- pure CPU, no network --- runs
ten-wide. A completion sentinel prevents a partial download being treated as
finished.

\paragraph{Reproducibility under change.} Extending the graph proved to be the
main source of defect. Adding a relation touches nine files, and on two
occasions a hardcoded list was not updated: the RDF was correct while the
export and the baselines silently described the previous state. The released
repository documents the checklist, and the pipeline now \emph{discovers}
plot references by scanning staged files for the identifying column rather than
consulting a list. We report this because the failure mode --- correct data,
stale derived artifacts --- is invisible to validation.

%======================================================================
\section{Validation: What SHACL Caught, and What It Did Not}
\label{sec:validation}

The gate comprises 22 node shapes with cardinality, datatype, range and closed
value-set constraints, plus SPARQL-based cross-record checks. It is a gate, not
a report: a site graph failing any violation-severity constraint is not
published.

\subsection{Three classes of defect only integration exposes}

\textbf{Dangling spatial references.} Drag-cloth bouts at 270 plots, and
phenology plants at 60 more, referenced plot IRIs that no product ever typed.
Each product was internally consistent; the defect existed only in the union.

\textbf{A referential gap in the source data.} Three phenophase observations at
one site concern a plant absent from NEON's own plant registry --- an
observation with no subject. This is a gap in published data, not in our
processing, and would be invisible to anyone reading either table alone.

\textbf{Cross-source assumptions that a second network breaks.} A shape written
when NEON was the only source required phenophase locations to be
\texttt{Plot}. USA-NPN observations occur at \texttt{ObservationSite}, correctly
so; the constraint produced 10{,}400 violations per site. The shape was wrong,
not the data. Adding a source, we conclude, requires re-reading every shape
that constrains a location or a class.

\subsection{Four errors SHACL could not catch}

All four passed the gate cleanly and were found only by inspecting value
distributions against expectation.

\begin{enumerate}
  \item NEON records pathogen assays, tick species identifications and DNA
        quality controls in one \texttt{testPathogenName} field. Counting all
        three as detections inflated tick-borne detections from 604 to 3{,}398.
        The tell was a ``pathogen'' with near-100\% positivity.
  \item Tick attachment is recorded as Y/N, not as a count. Parsing it as an
        integer silently discarded every value; the tell was a reported burden
        of zero on every animal.
  \item \texttt{individualCount} is 1 in all 231{,}892 tick records: these are
        individual specimens, not pools. A pooled-prevalence correction would
        have been wrong.
  \item Drag bouts finding nothing have no \texttt{sampleID} --- there is no
        sample to identify --- so requiring one discarded exactly the observed
        negatives that make vector presence a balanced task. The tell was a
        98\% presence rate where the raw data showed 43\%.
\end{enumerate}

Every one of these produced a well-formed graph. Structural validation proves
data well-formed; it does not prove it correct. We believe resource papers
should say so.

%======================================================================
\section{The Derived Benchmark}
\label{sec:benchmark}

\subsection{Deriving a temporal graph from RDF}

The derivation is a reusable pattern rather than a one-off script. Timestamped
events are read from the reified observation nodes, \emph{not} from the
materialised shortcut edges: \texttt{detectedIn} deduplicates 246 detections to
215 pathogen--host pairs, and a graph built from the shortcuts would lose the
temporal dimension entirely. Any SOSA-based observation graph faces the same
choice.

\begin{table}[t]
\centering
\caption{Dataset composition. Prevalence is the positive fraction among
evaluable events.}
\label{tab:dataset}
\small
\begin{tabular}{llrrr}
\toprule
\textbf{Relation} & \textbf{head $\rightarrow$ tail} & \textbf{Events} &
\textbf{Positives} & \textbf{Prev.} \\
\midrule
\texttt{detectedIn}  & pathogen $\rightarrow$ host    & 11{,}157  & 246 & 2.2\% \\
\texttt{detectedAt}  & pathogen $\rightarrow$ plot    & 16{,}244  & 358 & 2.2\% \\
\texttt{observedAt}  & vector $\rightarrow$ plot      & 13{,}425  & 8{,}020 & 60\% \\
\texttt{expressedAt} & phenophase $\rightarrow$ place & 578{,}636 & 263{,}223 & 45\% \\
\texttt{capturedAt}  & host $\rightarrow$ plot        & 106{,}364 & --- & context \\
\midrule
\textbf{Total} & & \textbf{725{,}826} & \textbf{271{,}847} & \\
\bottomrule
\end{tabular}
\end{table}

Two properties distinguish the dataset from existing benchmarks. It carries
347{,}615 \textbf{observed negatives} --- assays and bouts that were performed
and found nothing --- rather than only sampled ones. And label density spans
more than an order of magnitude across relations sharing nodes, geography and
time.

\subsection{Negative sampling changes the conclusions}

Under uniformly random negatives, three memorization baselines that measure
different things --- presence, frequency, recency --- returned \emph{identical}
scores to four decimal places, because an unseen pair is separable by presence
alone. We therefore ship 50\% historical negatives: destinations the source has
been linked to at another time. The baselines then separate, and the ordering
differs by relation. Reported numbers use the shipped configuration;
\texttt{--hist-neg-frac 0} reproduces the easier variant.

\subsection{Four relations, four optimal baselines}

\begin{table}[t]
\centering
\caption{Test MRR under the TGB one-vs-many protocol, 20 negatives per
positive. Each relation is compared against \emph{its own} random floor: the
detection relations have smaller negative pools, so their floors are higher.}
\label{tab:baselines}
\small
\begin{tabular}{lrrrlr}
\toprule
\textbf{Relation} & \textbf{$n$} & \textbf{Random} & \textbf{EdgeBank} &
\textbf{Best baseline} & \\
\midrule
\texttt{detectedIn}  &     37 & 0.246 & \textbf{0.112} & EdgeBank\,(365d)  & 0.177 \\
\texttt{detectedAt}  &     54 & 0.254 & 0.331 & EdgeBank\,($\infty$)      & 0.331 \\
\texttt{observedAt}  & 1{,}203 & 0.167 & 0.268 & Occurrence count         & 0.466 \\
\texttt{expressedAt} & 39{,}548 & 0.175 & 0.276 & Occurrence count        & 0.516 \\
\bottomrule
\end{tabular}
\end{table}

Table~\ref{tab:baselines} shows a gradient tracking how regular each underlying
process is. Phenology is seasonal clockwork and yields 0.516 to simple
occurrence counting. Vector presence is seasonal but noisier, 0.466. Pathogen
detection at a place is sparse and irregular, 0.331.

\textbf{And pathogen detection in an individual animal defeats memorization
entirely}, scoring 0.112 against its own 0.246 floor, with no baseline
recovering it. The mechanism is legible: a mouse is typically captured once, so
the true pathogen--host pair is novel, while the historical negatives are pairs
already observed. A model scoring by recurrence ranks the truth last. This
result survived five rebuilds and two negative-sampling schemes.

We are not aware of another temporal graph benchmark in which one relation
rewards unlimited memory, another rewards bounded recency, and a third makes
memory actively harmful. A method cannot succeed on all four by exploiting one
kind of structure.

%======================================================================
\section{Cross-Source Integration}
\label{sec:crosssource}

An ontology paper can assert interoperability. Demonstrating it requires a
second source.

The USA National Phenology Network's Nature's Notebook programme collects
phenophase status records to a protocol compatible with NEON's: is this
phenophase currently expressed on this plant, yes or no. Nine of NEON's ten
phenophase names appear verbatim in the USA-NPN vocabulary, because NEON adopted
USA-NPN's protocol terms --- and PPO was built to serve both. The join therefore
runs through the PPO terms already in our ontology. \textbf{No crosswalk was
written for the purpose.} A further 16 growth-form variants
(``Leaves (grasses)'', ``Breaking needle buds (conifers)'') were mapped
explicitly, raising retention from 56\% to 88\%; the remainder observe
phenophases NEON does not record and are excluded to keep the sources
comparable.

Restricting to stations within 2\,km of a NEON site yields 79 stations covering
38 of 39 sites at a \textbf{median separation of 230\,m} --- these are
deliberately co-located research installations, not incidental volunteer sites.

\begin{table}[t]
\centering
\caption{Mean day-of-year of first \emph{Breaking leaf buds} expression, by
network, at the six best-sampled co-located sites. Full results for all 27
paired sites are in the repository.}
\label{tab:agreement}
\small
\begin{tabular}{lrrr}
\toprule
\textbf{Site} & \textbf{NEON} & \textbf{USA-NPN} & \textbf{Difference} \\
\midrule
DELA & 150.3 & 150.0 & $-0.2$ \\
ORNL & 124.3 & 123.5 & $-0.8$ \\
TALL & 151.6 & 150.4 & $-1.2$ \\
SERC & 140.7 & 140.4 & $-0.3$ \\
UKFS & 132.0 & 131.4 & $-0.6$ \\
HARV & 147.7 & 148.1 & $+0.5$ \\
\midrule
\multicolumn{4}{l}{\emph{All 27 paired sites:} mean $-0.95$ d, median
absolute 1.59 d, 22/27 within 5 d} \\
\bottomrule
\end{tabular}
\end{table}

\textbf{The two networks agree.} Across 27 sites the mean difference is
$-0.95$ days and the median absolute difference 1.59 days, with no systematic
bias in either direction. Disagreement correlates with sample size
($r=-0.456$): the twelve sites with 570 or more paired observations all agree
within two days without exception, while the outliers have 3, 48 and 101
observations. One case resists that account --- JERC differs by 19.7 days on
331 observations --- and we report it unresolved; differing species composition
between the two sets of marked plants is the likeliest explanation.

This matters beyond validation of the mapping. USA-NPN states that its
coordinating office has limited ability to verify individual observer skill. At
co-located sites the volunteer aggregate nevertheless tracks trained
technicians to within about a day. That is evidence about citizen science data
quality, obtained as a by-product of demonstrating the alignment works.

We emphasise what made the question askable. Comparing two monitoring networks
requires a vocabulary in which their observations are the same kind of thing.
Every other analysis we performed on this graph could have been done with a
data frame; this one could not.

\begin{lstlisting}[caption={Both networks, one query, joined through shared PPO-aligned terms.}]
SELECT ?network (AVG(?doy) AS ?meanGreenup) WHERE {
  ?o  a neonepn:PhenophaseObservation ;
      neonepn:observedByNetwork   ?net ;
      neonepn:observedPhenophase  neonepn:BreakingLeafBuds ;
      neonepn:hasPhenophaseStatus neonepn:Expressed ;
      sosa:resultTime             ?ts .
  BIND(REPLACE(STR(?net), "^.*#", "") AS ?network)
} GROUP BY ?network
\end{lstlisting}

%======================================================================
\section{Ecological Findings}

The graph was built to answer integrative questions. We report the outcomes
briefly, including the negative ones, as evidence of what the resource
supports.

The pipeline \textbf{reproduces the known Lyme gradient} from raw products with
no prior knowledge of tick ecology: \emph{Borrelia burgdorferi} prevalence runs
9.5\% at a Wisconsin site against 0.18\% in Florida, with magnitudes matching
published nymphal infection rates. This is a validation rather than a
discovery, and serves as a regression test.

\textbf{Accumulated warmth predicts tick activity but not infection.}
Degree-days above 0\,\textdegree C in the 30 days before a drag bout predict
whether it finds ticks ($+0.00035$ per degree-day, 95\% CI
$[+0.00017, +0.00052]$, 8{,}275 bouts, 39 sites, site fixed effects, bootstrap
clustered on site) --- a 10.5 percentage-point difference between a
20\,\textdegree C and a 10\,\textdegree C month. Annual temperature does not
predict annual abundance, and no climate measure predicts infection. Questing
responds to warmth within weeks; infection reflects multi-year host dynamics.

\textbf{Three null results.} No dilution effect across three specifications, at
site and plot scale and on both infection and presence outcomes. No short-term
climate signal for infection. And spring timing proved \emph{unanswerable as
posed}: green-up does not predict tick onset, but a control regression revealed
onset tracking the \emph{sampling start date} instead --- the outcome variable
is censored by the field calendar. We report this because the analysis
diagnosed its own confound.

%======================================================================
\section{Availability, Licensing and Sustainability}

\begin{table}[h]
\centering
\small
\begin{tabular}{ll}
\toprule
Ontology, shapes, mappings & CC BY 4.0 \\
Pipeline and analysis code & MIT \\
Derived dataset            & CC BY 4.0 \\
Repository                 & \textit{[anonymised for review]} \\
Persistent archive         & \textit{[DOI at camera-ready]} \\
\bottomrule
\end{tabular}
\end{table}

NEON data moved from CC0 to CC BY 4.0 on 30 June 2026; USA-NPN data is CC BY
4.0. Both permit redistribution and derivative works with attribution. The
repository documents every transformation applied, as CC BY requires, and
carries the acknowledgement USA-NPN's attribution policy specifies.

\paragraph{Ontology evaluation.} The SHACL gate validates \emph{data}; it
says nothing about whether the model is well built. We assess the ontology
separately, and report what we could and could not run.

Competency-question coverage is reported in Section~\ref{sec:cq}: all seven
questions return results over a single site graph, and the SPARQL is
distributed with the resource. Structural metrics are computed locally
following the definitions OntoMetrics uses: 27 classes, 19 object and 17
datatype properties, 46 subclass axioms, maximum inheritance depth 3,
inheritance richness 1.70, attribute richness 0.63. Annotation coverage is
100\% for both \texttt{rdfs:label} and \texttt{rdfs:comment} across 87 terms
--- it was 78\% for comments until running the evaluation exposed 19
undocumented terms, which is a small but concrete argument for these
instruments. A reuse audit records one imported vocabulary and seven
referenced by IRI, with nine SKOS mappings to external terms.

We were unable to run OOPS!. Its hosted interface no longer presents an input
form, its REST endpoint returns HTTP 500, and the Docker image it now
recommends publishes no \texttt{linux/amd64} build, so it cannot run on the
x86 cluster available to us. We report this rather than omit the instrument,
since a reader attempting the same evaluation will meet the same obstacles.
FOOPS! assesses a resolvable URI rather than uploaded content and will be run
once the ontology is served from its namespace.

\paragraph{Reproducibility.} The pipeline runs end to end from five Slurm
submissions. Offline fixtures mirroring the real source schemas allow the full
path --- normalisation, RML materialisation, SHACL validation, export --- to be
exercised without network access or credentials.

\paragraph{Sustainability.} The ontology is versioned and published at a
resolvable namespace. Both upstream networks are ongoing national
infrastructure, so the pipeline re-runs against new releases without
modification; the site window configuration is the only thing that changes.

%======================================================================
\section{Conclusion}

NEON-EPN integrates six data products from two independently operated
observation networks into a validated knowledge graph, and derives from it a
temporal graph benchmark with properties we have not found elsewhere. The
artifact is available, documented and reproducible.

We would highlight two things for reuse. The RDF-to-temporal-graph derivation
--- read events from reified observations, never from materialised shortcut
edges --- applies to any SOSA-based observation graph, and we are not aware of
prior work in that direction. And the honest account of validation in
Section~\ref{sec:validation}: SHACL caught three classes of referential defect
that no single-product check could surface, and missed four modelling errors
that only distributional inspection revealed. Both facts are useful to anyone
building a comparable resource.

\paragraph{Limitations.} The two detection relations have 37 and 54 test
positives; conclusions about them are correspondingly weak, and we report them
separately rather than pooled. Baselines are memorization heuristics, and a
learned model has not yet been evaluated. The cross-source comparison uses an
approximate day-of-year and covers one phenophase.

%======================================================================
\section*{Resource Availability Statement}

Source code, ontology, SHACL shapes, RML mappings, the derived dataset and all
analysis scripts are available under the licences above. The anonymised
repository for review is cited in the submission system; the persistent
identifier will be supplied at camera-ready.

\bibliographystyle{plain}
\begin{thebibliography}{99}
\small

\bibitem{tgb2}
Gastinger, J., Huang, S., Galkin, M., Loghmani, E., Parviz, A., Poursafaei, F.,
Danovitch, J., Rossi, E., Koutis, I., Stuckenschmidt, H., Rabbany, R.,
Rabusseau, G.: TGB 2.0: A benchmark for learning on temporal knowledge graphs
and heterogeneous graphs. In: NeurIPS (2024)

\bibitem{edgebank}
Poursafaei, F., Huang, S., Pelrine, K., Rabbany, R.: Towards better evaluation
for dynamic link prediction. In: NeurIPS (2022)

\bibitem{ppo}
Stucky, B.J., Guralnick, R., Deck, J., Denny, E.G., Bolmgren, K., Walls, R.:
The Plant Phenology Ontology: A new informatics resource for large-scale
integration of plant phenology data. Frontiers in Plant Science \textbf{9},
517 (2018)

\bibitem{sosa}
Janowicz, K., Haller, A., Cox, S.J.D., Le Phuoc, D., Lefran\c{c}ois, M.: SOSA:
A lightweight ontology for sensors, observations, samples, and actuators.
Journal of Web Semantics \textbf{56}, 1--10 (2019)

\bibitem{neon}
Thibault, K.M., Laney, C.M., Yule, K.M., Franz, N.M., et al.: The US National
Ecological Observatory Network and the Global Biodiversity Framework. Journal
of Ecology and Environment \textbf{47} (2023)

\bibitem{neonticks}
Paull, S.H., et al.: Tick abundance, diversity and pathogen data collected by
the National Ecological Observatory Network. Ecology (2023)

\bibitem{neonderived}
Atkins, J.W., et al.: Recommendations for developing, documenting, and
distributing data products derived from NEON data. Ecosphere (2025)

\bibitem{qudtneon}
Porter, J.H., O'Brien, M., Frants, M., Earl, S., Martin, M., Laney, C.: Using
a units ontology to annotate pre-existing metadata. Scientific Data (2025)

\bibitem{morphkgc}
Arenas-Guerrero, J., Chaves-Fraga, D., Toledo, J., P\'erez, M.S.,
Corcho, O.: Morph-KGC: Scalable knowledge graph materialization with mapping
partitions. Semantic Web (2022)

\bibitem{shacl}
Knublauch, H., Kontokostas, D.: Shapes Constraint Language (SHACL). W3C
Recommendation (2017)

\bibitem{usanpn}
Rosemartin, A., Denny, E.G., Gerst, K.L., Marsh, R.L., Posthumus, E.E.,
Crimmins, T.M., Weltzin, J.F.: USA National Phenology Network observational
data documentation. USGS Open-File Report 2018--1060 (2018)

\bibitem{lot}
Poveda-Villal\'on, M., Fern\'andez-Izquierdo, A., Fern\'andez-L\'opez, M.,
Garc\'ia-Castro, R.: LOT: An industrial oriented ontology engineering
framework. Engineering Applications of Artificial Intelligence \textbf{111}
(2022)

\bibitem{ontopfas}
Di Pierro, D., Abrouk, L., Guyot, A., Symeonidou, D., Labadie, P.,
Lysaniuk, B.: OntoPFAS: Ontologie des PFAS et de leur exposition. In: PFIA
(2025)

\bibitem{oops}
Poveda-Villal\'on, M., G\'omez-P\'erez, A., Su\'arez-Figueroa, M.C.:
OOPS! (OntOlogy Pitfall Scanner!): An on-line tool for ontology evaluation.
International Journal on Semantic Web and Information Systems \textbf{10}(2)
(2014)

\bibitem{dilution}
Ostfeld, R.S., Keesing, F.: Biodiversity and disease risk: the case of Lyme
disease. Conservation Biology \textbf{14}(3), 722--728 (2000)

\end{thebibliography}

\end{document}
